% =====================================================================
%  Bölüm 12 — Tamlık Teoremi
% =====================================================================
\chapter{Tamlık Teoremi}\label{bol:tamlik}

\begin{bolumhedefleri}
\begin{itemize}
  \item Tamlık teoreminin ne söylediğini, zayıf ve güçlü biçimlerini
        ayırt etmeyi,
  \item Uzun bir metamantık ispatını küçük, bağımsız önsavlara bölmeyi,
  \item Kalmár'ın yöntemiyle doğruluk tablosunun her satırını bir
        türetime dönüştürmeyi,
  \item Atomları teker teker eleyerek önermeler mantığının tamlık
        teoremini \Lean{}'de eksiksiz ispatlamayı,
  \item Lindenbaum önsavı ve Henkin yöntemiyle birinci mertebe mantığın
        tamlık teoremine giden yolu görmeyi.
\end{itemize}
\end{bolumhedefleri}

Bölüm~11'de köprünün bir ayağını kurduk: Türetilebilen her şey
doğrudur. Bu bölümde öteki ayağı kuruyoruz:
\[
  \Gamma \models A \quad\Longrightarrow\quad \Gamma \vdash A.
\]
İki teorem birlikte, ispat sistemimizin iki değerli anlambilimi
\emph{tam olarak} yakaladığını söyler. Ne fazla şey türetir, ne eksik.
Geçerlilik ispatı, türetim ağacı üzerinde düz bir tümevarımdı. Tamlık
ise bambaşka bir iş ister: Elimizde yalnızca ``her değerlemede doğru''
bilgisi vardır ve bu bilgiden, hiç görmediğimiz bir türetim ağacını
\emph{inşa etmemiz} gerekir.

% ---------------------------------------------------------------------
\section{Tamlığın ifadesi}\label{sec:tamlik-ifade}

\begin{tanim}{Zayıf ve güçlü tamlık}{tamlik-turleri}
Bir ispat sistemi, bir anlambilime göre
\begin{itemize}
  \item \textbf{zayıf tamdır}, eğer her totoloji türetilebiliyorsa:
        $\models A \Rightarrow\; \vdash A$;
  \item \textbf{güçlü tamdır}, eğer her mantıksal sonuç
        türetilebiliyorsa: $\Gamma \models A \Rightarrow \Gamma \vdash
        A$. Burada $\Gamma$ sonsuz bir küme de olabilir.
\end{itemize}
\end{tanim}

\noindent Bizim \lc{Sonuc} ve \lc{Turetilir} tanımlarımızda $\Gamma$ bir
\emph{liste}, dolayısıyla sonludur. Sonlu $\Gamma$ için güçlü tamlık
zayıf tamlıktan kolayca çıkar: Varsayımları birer öncül olarak formülün
içine taşırız (\cref{sec:tamlik-teoremi}). Sonsuz $\Gamma$ için ise ek
bir fikir, \emph{tıkızlık} gerekir (\cref{sec:lindenbaum}).

\begin{tarihce}[Tarihçe: Tamlığın kısa tarihi]
Önermeler mantığının tamlığı ilk olarak Paul Bernays'in 1918 tarihli
(yayımlanmamış) habilitasyon tezinde ve bağımsız olarak Emil Post'un
1921'deki makalesinde ispatlandı \citep{post1921}. Post'un ispatı,
doğruluk tablolarının bugünkü biçimini de ilk kez sistemli olarak
kullanır. Bu bölümde izleyeceğimiz yöntem, Macar matematikçi László
Kalmár'ın 1935'te verdiği zarif ispattır \citep{kalmar1935}. Mendelson'ın
klasik ders kitabı da bu ispatı kullanır \citep{mendelson2015}. Birinci
mertebe mantığın tamlığını ise Kurt Gödel 1929'daki doktora teziyle
ispatladı \citep{godel1930}. Bugün ders kitaplarında okutulan ispat, Leon
Henkin'in 1949'daki ispatıdır \citep{henkin1949}.
\end{tarihce}

% ---------------------------------------------------------------------
\section{İspatın haritası}\label{sec:tamlik-harita}

Tamlık ispatı, kitabın şimdiye kadarki en uzun \Lean{} ispatı olacak.
Uzun bir ispatı yönetmenin yolu, onu her biri tek başına anlaşılabilen ve
tek başına derlenen \emph{önsavlara} bölmektir.
\Cref{fig:tamlik-harita}, ispatın beş aşamasını ve birbirine nasıl
bağlandığını gösteriyor.

\begin{figure}[htbp]
\centering
\begin{tikzpicture}[
    asama/.style={draw=#1, fill=#1!8, very thick, rounded corners=5pt,
                  text width=3.0cm, align=center, inner sep=5pt,
                  font=\small, minimum height=1.35cm},
    ok/.style={-{Stealth[length=2.6mm]}, very thick, black!45},
    no/.style={font=\baslikfont\scriptsize\bfseries, text=white,
               fill=#1, circle, inner sep=1.6pt},
  ]
  \node[asama=anarenk]     (a1) at (0,0)    {\textbf{Türetilmiş kurallar}\\{\scriptsize \lc{zay}, \lc{ucuncu\_hal}, \dots}};
  \node[asama=ornekframe]  (a2) at (4.2,0)  {\textbf{Kalmár önsavı}\\{\scriptsize her satır bir türetim}};
  \node[asama=ikincil]     (a3) at (8.4,0)  {\textbf{Atom eleme}\\{\scriptsize $p \oveya {\sim}p$ ile birleştir}};
  \node[asama=dikkatframe] (a4) at (8.4,-2.4) {\textbf{Zayıf tamlık}\\{\scriptsize $\models A \Rightarrow\; \vdash A$}};
  \node[asama=teoremframe] (a5) at (4.2,-2.4) {\textbf{Tamlık}\\{\scriptsize $\Gamma \models A \Rightarrow \Gamma \vdash A$}};
  \node[asama=arastirmaframe] (a6) at (0,-2.4) {\textbf{$\vdash \;=\; \models$}\\{\scriptsize geçerlilik ile birlikte}};
  \draw[ok] (a1) -- (a2); \draw[ok] (a2) -- (a3); \draw[ok] (a3) -- (a4);
  \draw[ok] (a4) -- node[above, font=\baslikfont\scriptsize]{\lc{kapat}} (a5);
  \draw[ok] (a5) -- (a6);
  \foreach \n/\c/\k in {a1/anarenk/1, a2/ornekframe/2, a3/ikincil/3, a4/dikkatframe/4, a5/teoremframe/5}
    \node[no=\c] at (\n.north west) {\k};
\end{tikzpicture}
\caption{Tamlık ispatının haritası. Her kutu, \lc{Bolum12.lean}
dosyasında bir ya da birkaç önsava karşılık gelir.}
\label{fig:tamlik-harita}
\end{figure}

\noindent Ana fikir şudur. Bir $A$ totolojisinin doğruluk tablosunun her
satırı, atomlara birer değer verir. Kalmár, \emph{her satırın kendisinin
bir türetime dönüştürülebildiğini} gözlemledi. Satırda doğru olan
atomları, yanlış olanların ise değillerini varsayarsak, bu varsayımlardan
$A$ türetilebilir. Sonra satırları ikişer ikişer birleştirerek
varsayımlardan kurtuluruz.

\subsection{Türetilmiş kurallar}

\Lean{}'deki \lc{Turetilir} yüklemi yalnızca on bir temel kurala sahip.
Uzun türetimler yazarken aynı küçük kalıpları tekrar tekrar kullanırız.
Matematikte bir önsavı bir kez ispatlayıp sonra kullanırız; aynı şeyi
ispat sistemimizde de yapabiliriz. Bir \emph{türetilmiş kural},
temel kurallardan ispatlanan ve sonra temel kural gibi kullanılan bir
teoremdir.

En sık kullanacağımız iki araç, bağlamın başındaki varsayımı kullanmak ve
bağlamın başına yeni bir varsayım eklemektir:

\leandosya[Temel araçlar]{Bolum12}{temel_araclar}

\noindent Klasik mantığın kalbi olan üçüncü hâlin olanaksızlığını
Bölüm~11'in alıştırmalarında boş bağlamda türetmiştik (\lc{a11\_7}).
Şimdi onu \emph{her} bağlamda ve \emph{her} formül için türetiyor, ondan
da \emph{durum ayrımı} kuralını çıkarıyoruz:

\leandosya[Üçüncü hâl ve durum ayrımı]{Bolum12}{ucuncu_hal}

\begin{center}
\begin{prooftree}
  \hypo{B, \Gamma \vdash A}
  \hypo{{\sim}B, \Gamma \vdash A}
  \infer2[Durum]{\Gamma \vdash A}
\end{prooftree}
\end{center}

\noindent Bu kural, ispatın son aşamasında atomları elemek için
kullanılacak. Kalmár önsavı içinse her bağlacın doğruluk tablosunu
taklit eden kurallara ihtiyacımız var (\cref{tab:turetilmis-kurallar}).

\leandosya[Bağlaçlar için türetilmiş kurallar]{Bolum12}{turetilmis_kurallar}

\begin{table}[htbp]
\centering
\small
\renewcommand{\arraystretch}{1.5}
\rowcolors{2}{tablosatir}{white}
\begin{tabular}{>{\raggedright\arraybackslash}p{3.0cm}
                >{\centering\arraybackslash}p{4.4cm}
                >{\raggedright\arraybackslash}p{5.0cm}}
\rowcolor{tabloust}
\tablobaslik{Lean adı} & \tablobaslik{Kural} & \tablobaslik{Doğruluk tablosundaki karşılığı} \\
\lc{ve\_yanlis1}    & ${\sim}A \;\vdash\; {\sim}(A \ove B)$          & $A$ yanlışsa $A \ove B$ yanlış \\
\lc{ve\_yanlis2}    & ${\sim}B \;\vdash\; {\sim}(A \ove B)$          & $B$ yanlışsa $A \ove B$ yanlış \\
\lc{veya\_yanlis}   & ${\sim}A,\, {\sim}B \;\vdash\; {\sim}(A \oveya B)$ & ikisi de yanlışsa $A \oveya B$ yanlış \\
\lc{ise\_on\_yanlis} & ${\sim}A \;\vdash\; A \oise B$                 & öncül yanlışsa $A \oise B$ doğru \\
\lc{ise\_son\_dogru} & $B \;\vdash\; A \oise B$                      & sonuç doğruysa $A \oise B$ doğru \\
\lc{ise\_yanlis}    & $A,\, {\sim}B \;\vdash\; {\sim}(A \oise B)$    & tek yanlış satır \\
\end{tabular}
\caption{Kalmár önsavı için gereken türetilmiş kurallar. Doğru durumlar
($A, B \vdash A \ove B$ gibi) zaten temel kurallardır.}
\label{tab:turetilmis-kurallar}
\end{table}

\begin{onemlifikir}[Önemli Fikir: İspat sistemi içinde ispat sistemi]
Türetilmiş kurallar, \Lean{}'in kendi ispatlarında yaptığımız şeyin nesne
dilindeki yansımasıdır. \lc{ve\_yanlis1} bir kez ispatlandıktan sonra
\lc{Turetilir}'in on ikinci kuralıymış gibi kullanılabilir. Sistemi
değiştirmeyiz; yalnızca onun hakkında bir teorem ispatlarız. Bu, ispat
sistemini küçük tutup (geçerlilik ispatı kısa kalsın diye) kullanımını
zenginleştirmenin standart yoludur.
\end{onemlifikir}

% ---------------------------------------------------------------------
\section{Kalmár yöntemi}\label{sec:kalmar}

\subsection{İşaretli atomlar}

Bir $v$ değerlemesi ve bir $p_n$ atomu için \emph{işaretli atomu}
\[
  p_n^v \;=\;
  \begin{cases}
    p_n        & v(n) = \mathsf{true} \text{ ise},\\
    {\sim}p_n  & v(n) = \mathsf{false} \text{ ise}
  \end{cases}
\]
diye tanımlarız. $p_n^v$, $v$ altında her zaman doğrudur. İlk $k$ atomun
işaretli hâllerinden oluşan bağlama $\Delta_v^k = [p_{k-1}^v, \dots,
p_1^v, p_0^v]$ diyelim. Bu bağlam, doğruluk tablosunun bir satırının
nesne dilindeki ifadesidir.

\leandosya[İşaretli atomlar]{Bolum12}{isaretli}

\noindent \lc{atomBaglami}'nın sayı üzerinde özyinelemeyle tanımlanması,
ileride atomları teker teker elemeyi kolaylaştıracak: \lc{atomBaglami v
(k + 1)} bağlamının başında her zaman $p_k^v$ durur. İki basit özellik
daha ispatlıyoruz. İşaretli atomun açık biçimi ve bağlamda bulunması:

\leandosya[İşaretli atomların özellikleri]{Bolum12}{isaretli_ozellik}

\subsection{Kalmár önsavı}

\begin{teorem}{Kalmár önsavı}{kalmar}
$A$'nın bütün atomları $p_0, \dots, p_{k-1}$ arasındaysa, her $v$ için
\[
  \Delta_v^k \vdash A^v,
  \qquad\text{burada}\qquad
  A^v =
  \begin{cases}
    A        & \lBrack A \rBrack_v = \mathsf{true} \text{ ise},\\
    {\sim}A  & \lBrack A \rBrack_v = \mathsf{false} \text{ ise}.
  \end{cases}
\]
\end{teorem}

\noindent Önsav, \emph{doğru olanı da yanlış olanı da} türetebildiğimizi
söylüyor. Bu güçlendirme, tümevarımın işlemesi için şarttır. $A \oise B$
yanlış olduğunda ${\sim}(A \oise B)$'yi türetmek için $A$'yı ve
${\sim}B$'yi türetebiliyor olmamız gerekir. Yani tümevarım hipotezinin
hem doğru hem yanlış yüzüne ihtiyaç vardır (Bölüm~7'deki ``önermeyi
güçlendirme'' tekniği). \Cref{fig:kalmar-tablo}, önsavın $p \oise q$ için
ne söylediğini satır satır gösteriyor.

\begin{figure}[htbp]
\centering
\renewcommand{\arraystretch}{1.6}
\setlength{\tabcolsep}{9pt}
\begin{tabular}{cc c >{\centering\arraybackslash}p{3.0cm} c >{\raggedright\arraybackslash}p{3.3cm}}
\rowcolor{tabloust}
\tablobaslik{$p$} & \tablobaslik{$q$} & \tablobaslik{$p \oise q$} &
\tablobaslik{Bağlam $\Delta_v$} & & \tablobaslik{Türetilen ($A^v$)} \\
\cellcolor{ikincil!15}\textsf{D} & \cellcolor{ikincil!15}\textsf{D} & \textsf{D} &
  $q,\; p$ & $\vdash$ & $p \oise q$ \quad{\scriptsize\textcolor{leantac}{\lc{ise\_son\_dogru}}} \\
\cellcolor{ikincil!15}\textsf{D} & \cellcolor{vurgu!15}\textsf{Y} & \cellcolor{vurgu!15}\textsf{Y} &
  ${\sim}q,\; p$ & $\vdash$ & \textcolor{teoremframe}{${\sim}(p \oise q)$} \quad{\scriptsize\textcolor{leantac}{\lc{ise\_yanlis}}} \\
\cellcolor{vurgu!15}\textsf{Y} & \cellcolor{ikincil!15}\textsf{D} & \textsf{D} &
  $q,\; {\sim}p$ & $\vdash$ & $p \oise q$ \quad{\scriptsize\textcolor{leantac}{\lc{ise\_on\_yanlis}}} \\
\cellcolor{vurgu!15}\textsf{Y} & \cellcolor{vurgu!15}\textsf{Y} & \textsf{D} &
  ${\sim}q,\; {\sim}p$ & $\vdash$ & $p \oise q$ \quad{\scriptsize\textcolor{leantac}{\lc{ise\_on\_yanlis}}} \\
\end{tabular}
\caption{Kalmár önsavı $p \oise q$ için. Doğruluk tablosunun her satırı
bir türetime dönüşür: Bağlam satırdaki atom değerlerini, sağ taraf ise
formülün o satırdaki değerini kodlar. Son sütun, satırın hangi türetilmiş
kuralla elde edildiğini gösterir.}
\label{fig:kalmar-tablo}
\end{figure}

\noindent Önsavı \Lean{}'de ifade etmek için ``$b$ değerine uygun
biçimde türetir'' anlamına gelen küçük bir yüklem tanımlıyoruz. $A^v$'yi
bir \leankw{if} ifadesiyle tanımlamak yerine iki yönü bir \lc{∧} ile
yazmak, ispatlarda \leankw{if} ifadelerini açma zahmetinden kurtarır:

\leandosya[Kalmár yüklemi]{Bolum12}{kalmar_tanim}

\noindent \leankw{abbrev} kullanmamızın sebebi, \lc{Kalmar}'ın \lc{∧}
olarak açılabilmesidir. Böylece \lc{⟨\_, \_⟩} ile kurulabilir ve
\lc{hA.1}, \lc{hA.2} ile parçalanabilir. Şimdi her bağlaç için ayrı bir
önsav ispatlıyoruz. Her önsav, o bağlacın doğruluk tablosunu satır
satır ele alır:

\leandosya[Bağlaç adımları]{Bolum12}{kalmar_baglaclar}

\noindent Bu önsavlarda \lc{a} ve \lc{b}, alt formüllerin doğruluk
değerleridir. \leantac{cases a with} her satırı ayrı bir hedefe
dönüştürür. Olanaksız yönler (örneğin \lc{(false \&\& b) = true}) tek bir
\leantac{simp at h} ile kapanır. Geri kalan her hedef, ya bir temel
kuralla ya da \cref{tab:turetilmis-kurallar}'deki bir türetilmiş
kuralla tek satırda biter.

Önsavın kendisi artık yalnızca bir yapısal tümevarımdır. Her durum, ilgili
bağlaç önsavına yönlendirilir:

\leandosya[Kalmár önsavı]{Bolum12}{kalmar}

\begin{dikkat}[Dikkat: Tanımsal eşitlik iş başında]
\lc{kalmar\_ve}'nin sonucu \lc{Kalmar Δ (a \&\& b) (A {\sembolkodfont ⋀} B)} biçimindedir.
Hedef ise \lc{Kalmar Δ (deger v (A {\sembolkodfont ⋀} B)) (A {\sembolkodfont ⋀} B)}'dir. \Lean{} bu ikisini
\emph{tanımsal olarak} eşit görür, çünkü \lc{deger}'in tanımı açıldığında
ikisi aynı terim olur. Bu sayede \leantac{exact} hiçbir yeniden yazma
olmadan çalışır. \lc{deger}'i bir tablo yerine özyinelemeli bir fonksiyon
olarak tanımlamanın bir kazancı da budur.
\end{dikkat}

% ---------------------------------------------------------------------
\section{Atomların elenmesi}\label{sec:atom-eleme}

$A$ bir totolojiyse, Kalmár önsavı \emph{her} $v$ için $\Delta_v^k
\vdash A$ verir; çünkü $A^v$ her zaman $A$'nın kendisidir. Elimizde
$2^k$ tane türetim var, her satır için bir tane. Bunları birleştirip
bağlamı boşaltmak istiyoruz.

Fikir şudur: Yalnızca $p_{k-1}$'in değerinde farklılaşan iki satırı
alalım. Bu satırların bağlamları
\[
  p_{k-1},\, \Delta \vdash A
  \qquad\text{ve}\qquad
  {\sim}p_{k-1},\, \Delta \vdash A
\]
biçimindedir; durum ayrımı kuralı $\Delta \vdash A$ verir. Böylece
$p_{k-1}$'i bağlamdan çıkarmış oluruz. Bunu $k$ kez tekrarlarız
(\cref{fig:atom-eleme}).

\begin{figure}[htbp]
\centering
\begin{tikzpicture}[
    yaprak/.style={draw=ornekframe, fill=ornekbg, thick, rounded corners=3pt,
                   font=\small, inner sep=4pt, minimum width=2.6cm},
    ara/.style={draw=ikincil, fill=fikirbg, thick, rounded corners=3pt,
                font=\small, inner sep=4pt, minimum width=2.4cm},
    kok/.style={draw=teoremframe, fill=teorembg, very thick, rounded corners=4pt,
                font=\small\bfseries, inner sep=5pt},
    ok/.style={-{Stealth[length=2.2mm]}, thick, black!45},
    etiket/.style={font=\baslikfont\scriptsize\bfseries, text=ikincil},
  ]
  \node[yaprak] (l1) at (-5.1,0) {$q,\; p \vdash A$};
  \node[yaprak] (l2) at (-1.7,0) {${\sim}q,\; p \vdash A$};
  \node[yaprak] (l3) at (1.7,0)  {$q,\; {\sim}p \vdash A$};
  \node[yaprak] (l4) at (5.1,0)  {${\sim}q,\; {\sim}p \vdash A$};
  \node[ara] (m1) at (-3.4,-1.7) {$p \vdash A$};
  \node[ara] (m2) at (3.4,-1.7)  {${\sim}p \vdash A$};
  \node[kok] (r) at (0,-3.3) {$\vdash A$};
  \draw[ok] (l1) -- (m1); \draw[ok] (l2) -- (m1);
  \draw[ok] (l3) -- (m2); \draw[ok] (l4) -- (m2);
  \draw[ok] (m1) -- (r);  \draw[ok] (m2) -- (r);
  \node[etiket] at (-3.4,-0.85) {$q \oveya {\sim}q$};
  \node[etiket] at (3.4,-0.85)  {$q \oveya {\sim}q$};
  \node[etiket] at (0,-2.35)    {$p \oveya {\sim}p$};
  \node[font=\baslikfont\scriptsize, text=ornekframe, anchor=west] at (6.6,0)
       {\begin{tabular}{@{}l@{}}Kalmár:\\ $2^k$ satır\end{tabular}};
  \node[font=\baslikfont\scriptsize, text=ikincil, anchor=west] at (6.6,-1.7)
       {\begin{tabular}{@{}l@{}}\lc{atom\_ele}\\(1. adım)\end{tabular}};
  \node[font=\baslikfont\scriptsize, text=teoremframe, anchor=west] at (6.6,-3.3)
       {\begin{tabular}{@{}l@{}}\lc{atom\_ele}\\(2. adım)\end{tabular}};
\end{tikzpicture}
\caption{İki atomlu bir $A$ totolojisi için atomların elenmesi.
Yapraklar Kalmár önsavından gelir. Her birleştirme, bir atom için durum
ayrımıdır. $k$ atom için ağaç $k$ katlıdır ve $2^k$ yaprağı vardır.}
\label{fig:atom-eleme}
\end{figure}

\noindent Bunu \Lean{}'de yapmak için, bir $v$ değerlemesinden yalnızca
$j$'inci atomu değiştirilmiş yeni değerlemeler üretebilmeliyiz. Ayrıca
$j$'den küçük atomlarda anlaşan değerlemelerin aynı atom bağlamını
verdiğini bilmeliyiz:

\leandosya[Değerlemeyi güncellemek]{Bolum12}{ata}

\noindent Artık eleme adımını yazabiliriz. Hipotez, $j+1$ atomlu
\emph{her} satırın $A$'yı türettiğini söylüyor. $v$'yi verilmiş kabul
edip $j$'inci atomu bir kez \lc{true}, bir kez \lc{false} yapıyoruz.
İki türetimi durum ayrımıyla birleştirince $j$ atomlu satır elde ediliyor:

\leandosya[Atom eleme]{Bolum12}{atom_ele}

\noindent \lc{atomlari\_ele}, \lc{atom\_ele}'yi $k$ kez uygular. Sayı
üzerinde yapısal özyinelemeyle tanımlanmış bir \emph{teoremdir}:
\lc{k + 1} durumu, \lc{k} durumunu çağırır. $k = 0$ olduğunda atom
bağlamı boştur, dolayısıyla herhangi bir değerleme (burada hep yanlış
olanı) işimizi görür.

% ---------------------------------------------------------------------
\section{Tamlık teoremi}\label{sec:tamlik-teoremi}

Parçaları birleştirelim. $A$'nın atomları $p_0, \dots, p_m$
arasındadır; burada $m$, \lc{enBuyukAtom A} değeridir. Dolayısıyla
$k = m + 1$ yeterlidir.

\begin{teorem}{Zayıf tamlık}{zayif-tamlik}
Her $A$ totolojisi türetilebilir: $\models A \Rightarrow\; \vdash A$.
\end{teorem}

\leandosya[Zayıf tamlık]{Bolum12}{zayif_tamlik}

\noindent Varsayımlı hâl için, varsayımları öncül olarak formülün içine
taşıyan bir \lc{kapat} fonksiyonu tanımlıyoruz. Sonra iki şey
gösteriyoruz: Anlamsal sonuç ilişkisi bu taşımayla korunur, türetim ise
geri alınabilir. İkincisi, Bölüm~5'teki \emph{çıkarım teoreminin} ($\to$G
kuralının) tersidir:

\leandosya[Varsayımları içeri taşımak]{Bolum12}{kapat}

\begin{teorem}{Tamlık (önermeler mantığı)}{tamlik}
Her sonlu $\Gamma$ ve her $A$ için, $\Gamma \models A$ ise $\Gamma
\vdash A$. Geçerlilikle birlikte:
\[
  \Gamma \vdash A \quad\Longleftrightarrow\quad \Gamma \models A .
\]
\end{teorem}

\leandosya[Tamlık teoremi]{Bolum12}{tamlik}

\begin{figure}[htbp]
\centering
\begin{tikzpicture}[
    kutu/.style={draw=#1, fill=#1!8, very thick, rounded corners=6pt,
                 text width=3.6cm, align=center, inner sep=8pt},
  ]
  \node[kutu=anarenk] (s) at (0,0)
       {\baslikfont\bfseries Sözdizim\\[2pt]\normalfont $\Gamma \vdash A$\\[2pt]
        {\scriptsize türetim ağaçları}\\{\scriptsize \lc{Turetilir}}};
  \node[kutu=ikincil] (a) at (8,0)
       {\baslikfont\bfseries Anlambilim\\[2pt]\normalfont $\Gamma \models A$\\[2pt]
        {\scriptsize değerlemeler}\\{\scriptsize \lc{Sonuc}}};
  \draw[-{Stealth[length=3.2mm]}, very thick, teoremframe]
       ([yshift=6pt]s.east) -- node[above, font=\baslikfont\small\bfseries]{geçerlilik}
       node[below, font=\scriptsize]{Bölüm 11} ([yshift=6pt]a.west);
  \draw[-{Stealth[length=3.2mm]}, very thick, dikkatframe]
       ([yshift=-18pt]a.west) -- node[below, font=\baslikfont\small\bfseries]{tamlık}
       node[above, font=\scriptsize]{Bölüm 12} ([yshift=-18pt]s.east);
\end{tikzpicture}
\caption{Önermeler mantığının iki yüzü. Geçerlilik ve tamlık birlikte,
türetilebilirlik ile mantıksal sonucun \emph{aynı} ilişki olduğunu
söyler: \lc{turetilir\_iff\_sonuc}.}
\label{fig:iki-yuz}
\end{figure}

\begin{onemlifikir}[Önemli Fikir: Tamlık bir ispat \emph{üretir}]
Tamlık ispatımız yapıcıdır: \lc{tamlik\_bos}, verilen bir totoloji
için gerçek bir türetim ağacı inşa eder. Bu ağaç büyüktür. $k$ atomlu bir
formül için $2^k$ Kalmár türetimini birleştirir. Ama vardır ve
mekaniktir. Bu yüzden önermeler mantığında ``türetilebilir mi?''
sorusu, ``doğruluk tablosu her satırda doğru mu?'' sorusuna indirgenir.
Bölüm~10'daki \lc{totolojiMi} fonksiyonu, aslında bir \emph{ispat arama}
algoritmasıdır.
\end{onemlifikir}

% ---------------------------------------------------------------------
\section{Sonuçlar}\label{sec:tamlik-sonuclar}

Tamlığın ilk getirisi pratiktir: Artık bir türetimi elle kurmak zorunda
değiliz. Bölüm~11'de Peirce yasasını on adımlık bir türetimle elde
etmiştik. Şimdi doğruluk tablosunu kontrol etmek yeterli. İkinci getiri
kuramsaldır: Tutarlılık ile sağlanabilirlik aynı şeydir.

\leandosya[Tamlığın sonuçları]{Bolum12}{sonuclar}

\begin{teorem}{Model varlığı}{model-varligi}
$\Gamma \nvdash \bot$ ise $\Gamma$'nın bütün formüllerini doğrulayan bir
$v$ değerlemesi vardır.
\end{teorem}

\noindent Bu teoremin ispatında \lc{Classical.byContradiction}
kullandık, çünkü değerlemenin \emph{varlığını} çelişki yoluyla
gösteriyoruz. Aslında burada da yapıcı bir ispat mümkündür:
Formülün atomları sonlu olduğu için doğruluk tablosunun satırlarını tek
tek deneyebiliriz (\lc{tumDegerlemeler}). Bu ispat \lc{a12\_8}'in
ötesinde bir alıştırma olarak bırakılmıştır.

\begin{dikkat}[Dikkat: Hesaplama karmaşıklığı]
\lc{tamlik\_bos}'un ürettiği türetim, atom sayısında üsteldir. Daha
kısa türetimler var mı? Bu, \emph{ispat karmaşıklığı} alanının temel
sorusudur. Cook ve Reckhow, her totolojinin polinom boyutlu bir türetimi
olan bir ispat sistemi bulunmasının $\mathrm{NP} = \mathrm{coNP}$'ye denk
olduğunu gösterdi \citep{cookreckhow1979}. Bu, açık bir problemdir. Belirli
sistemler için alt sınırlar bilinir: Örneğin güvercin yuvası ilkesinin
\emph{çözünürlük} (resolution) türetimleri üstel uzunluktadır
\citep{haken1985}. Yani tamlık ``bir türetim vardır'' der; ``kısa bir
türetim vardır'' demez.
\end{dikkat}

% ---------------------------------------------------------------------
\section{Maksimal tutarlı kümeler ve Lindenbaum önsavı}\label{sec:lindenbaum}

Kalmár yöntemi, sonlu sayıda atom ve sonlu bir $\Gamma$ için çalışır.
$\Gamma$ sonsuz olduğunda, ya da birinci mertebe mantığa geçtiğimizde,
doğruluk tablosu kullanılamaz. Bu durumlarda kullanılan genel yöntem,
modeli \emph{sözdizimin kendisinden} inşa etmektir. Bu kesimin
ispatlarını kâğıt üzerinde veriyoruz. Birinci mertebe mantık için
\Lean{} formalizasyonu Cilt~II'nin ilk bölümünün konusudur.

\begin{tanim}{Maksimal tutarlı küme}{maksimal-tutarli}
Bir $\Delta$ formül kümesi
\begin{itemize}
  \item \textbf{tutarlıdır}, eğer $\Delta \nvdash \bot$ ise;
  \item \textbf{maksimal tutarlıdır}, eğer tutarlıysa ve kendisine
        eklenen her yeni formül onu tutarsız yapıyorsa.
\end{itemize}
\end{tanim}

\noindent Maksimal tutarlı bir $\Delta$, ``her soruya kesin bir cevap
veren'' bir kümedir. Her $A$ için ya $A \in \Delta$ ya da ${\sim}A \in
\Delta$'dır, ikisi birden asla. Üstelik $\Delta$ türetim altında
kapalıdır ve bağlaçlarla uyumludur:
\[
  A \oise B \in \Delta \iff (A \notin \Delta \text{ ya da } B \in \Delta),
  \qquad
  A \ove B \in \Delta \iff (A \in \Delta \text{ ve } B \in \Delta).
\]
Bu yüzden $v_\Delta(n) = \mathsf{true} \iff p_n \in \Delta$ diye
tanımlanan \emph{kanonik değerleme}, yapısal tümevarımla şunu sağlar:
$\lBrack A \rBrack_{v_\Delta} = \mathsf{true} \iff A \in \Delta$
(\emph{doğruluk önsavı}).

\begin{teorem}{Lindenbaum önsavı}{lindenbaum}
Her tutarlı $\Gamma$ kümesi, maksimal tutarlı bir $\Delta \supseteq
\Gamma$ kümesine genişletilebilir.
\end{teorem}

\noindent \emph{İspat taslağı.} Formüller sayılabilir olduğu için onları
$\varphi_0, \varphi_1, \varphi_2, \dots$ diye sıralayabiliriz
(\lc{Formul} için böyle bir sıralama, Bölüm~8'deki sayılabilirlik
tekniğiyle kurulur). Sonra adım adım genişletiriz
(\cref{fig:lindenbaum}):
\[
  \Delta_0 = \Gamma, \qquad
  \Delta_{n+1} =
  \begin{cases}
    \Delta_n \cup \{\varphi_n\} & \text{bu küme tutarlıysa},\\
    \Delta_n \cup \{{\sim}\varphi_n\} & \text{değilse},
  \end{cases}
  \qquad
  \Delta = \bigcup_n \Delta_n .
\]
Her adım tutarlılığı korur. Çünkü hem $\Delta_n, \varphi_n \vdash \bot$
hem $\Delta_n, {\sim}\varphi_n \vdash \bot$ olsaydı, durum ayrımı
$\Delta_n \vdash \bot$ verirdi. Birleşim de tutarlıdır, çünkü bir
türetim yalnızca sonlu sayıda varsayım kullanır ve bunlar bir
$\Delta_n$'nin içindedir. \hfill\textcolor{teoremframe}{\rule{0.55em}{0.55em}}

\begin{figure}[htbp]
\centering
\begin{tikzpicture}[
    kume/.style={draw=#1, fill=#1!7, thick, rounded corners=6pt},
  ]
  \draw[kume=teoremframe] (-0.3,-1.25) rectangle (11.0,1.25);
  \draw[kume=ikincil]     (-0.15,-1.0) rectangle (7.6,1.0);
  \draw[kume=ornekframe]  (0.0,-0.8) rectangle (5.1,0.8);
  \draw[kume=anarenk]     (0.15,-0.6) rectangle (2.6,0.6);
  \node[font=\small] at (1.35,0) {$\Delta_0 = \Gamma$};
  \node[font=\small] at (3.85,0.2) {$\Delta_1$};
  \node[font=\small] at (6.35,0.2) {$\Delta_2$\;\;$\cdots$};
  \node[font=\small\bfseries, text=teoremframe] at (9.3,0) {$\Delta = \bigcup_n \Delta_n$};
  \node[font=\baslikfont\scriptsize, text=ornekframe] at (3.85,-0.45) {$+\,\varphi_0$ ya da ${\sim}\varphi_0$};
  \node[font=\baslikfont\scriptsize, text=ikincil] at (6.35,-0.55) {$+\,\varphi_1$ ya da ${\sim}\varphi_1$};
  \node[font=\baslikfont\scriptsize, text=teoremframe, anchor=north] at (9.3,-0.35) {maksimal tutarlı};
\end{tikzpicture}
\caption{Lindenbaum yapısı. Her adımda sıradaki formül ya da değili
eklenir; tutarlılık her adımda korunur. Sonsuz birleşim, her formül
hakkında karar vermiş maksimal tutarlı bir kümedir.}
\label{fig:lindenbaum}
\end{figure}

\noindent Lindenbaum önsavı ile kanonik değerleme birleşince, model
varlığı teoremi \emph{sonsuz} $\Gamma$ için de elde edilir: Tutarlı
$\Gamma$'yı maksimal tutarlı $\Delta$'ya genişletiriz ve $v_\Delta$,
$\Gamma$'nın her formülünü doğrular. Buradan güçlü tamlık çıkar:
$\Gamma \nvdash A$ ise $\Gamma \cup \{{\sim}A\}$ tutarlıdır, bir modeli
vardır ve bu model $\Gamma \nvDash A$'ya tanıklık eder.

\begin{teorem}{Tıkızlık}{tikizlik}
Bir $\Gamma$ kümesinin her sonlu alt kümesi sağlanabilirse, $\Gamma$'nın
kendisi de sağlanabilirdir.
\end{teorem}

\noindent \emph{İspat.} $\Gamma$ sağlanamasaydı, güçlü tamlıkla $\Gamma
\vdash \bot$ olurdu. Bu türetim sonlu sayıda varsayım kullanır; o sonlu
alt küme geçerlilik teoremiyle sağlanamaz olurdu. \hfill\textcolor{teoremframe}{\rule{0.55em}{0.55em}}

\begin{dikkat}[Dikkat: Lindenbaum ve seçim]
Formüller sayılabilir olduğu için yukarıdaki yapı seçim aksiyomu
gerektirmez. Ama ``$\Delta_n \cup \{\varphi_n\}$ tutarlı mı?'' sorusuna
her adımda cevap verebilmemiz gerekir. Önermeler mantığında bu karar
verilebilirdir (tamlık ile doğruluk tablosuna indirgenir). Birinci
mertebe mantıkta değildir (\ciltiki{4}). Bu yüzden \Lean{}'de Lindenbaum
yapısı genellikle \leankw{Classical} ile, ya da sayılamaz diller için
Zorn önsavıyla yapılır. Mathlib'deki \lc{zorn\_subset} buna hazır bir
araçtır.
\end{dikkat}

% ---------------------------------------------------------------------
\section{Birinci mertebe mantık: Henkin yöntemi}\label{sec:henkin}

Birinci mertebe mantıkta bir model, bir \emph{tanım evreni} ve bu evren
üzerinde yüklemlerin yorumlarından oluşur. Model varlığını göstermek
için Henkin'in fikri, evreni \emph{terimlerin kendisinden} kurmaktır.
Ancak bir sorun vardır. $\exists x\, \varphi(x) \in \Delta$ olabilir,
ama hiçbir $t$ terimi için $\varphi(t) \in \Delta$ olmayabilir. Terim
modelinde varoluş iddiasının bir tanığı yoktur.

Henkin'in çözümü, dile her varoluş formülü için yeni bir \emph{tanık
sabiti} eklemektir: $c_{\exists x\varphi}$. Bu sabitle birlikte
\[
  \exists x\, \varphi(x) \;\oise\; \varphi(c_{\exists x\varphi})
\]
\emph{Henkin aksiyomu} da kümeye katılır. Bu aksiyomların eklenmesi
tutarlılığı korur. Sonra Lindenbaum önsavı uygulanır ve terim modeli
tanımlanır (\cref{fig:henkin}).

\begin{figure}[htbp]
\centering
\begin{tikzpicture}[
    adim/.style={draw=#1, fill=#1!8, very thick, rounded corners=5pt,
                 text width=2.55cm, align=center, inner sep=5pt,
                 font=\small, minimum height=1.6cm},
    ok/.style={-{Stealth[length=2.6mm]}, very thick, black!45},
  ]
  \node[adim=anarenk]     (a) at (0,0)     {tutarlı\\$\Gamma$};
  \node[adim=ornekframe]  (b) at (3.3,0)   {Henkin\\genişlemesi\\{\scriptsize tanık sabitleri}};
  \node[adim=ikincil]     (c) at (6.6,0)   {maksimal\\tutarlı $\Delta$\\{\scriptsize Lindenbaum}};
  \node[adim=teoremframe] (d) at (9.9,0)   {terim\\modeli $\mathcal{M}_\Delta$\\{\scriptsize $\mathcal{M}_\Delta \models \Gamma$}};
  \draw[ok] (a) -- (b); \draw[ok] (b) -- (c); \draw[ok] (c) -- (d);
\end{tikzpicture}
\caption{Henkin'in tamlık ispatının dört adımı. Evren, terimlerden oluşur
(eşitlik varsa eşit terimlerin denklik sınıflarından). $\Delta$'nın
kapalılık özellikleri, terim modelinin $\Delta$'yı, dolayısıyla
$\Gamma$'yı doğrulamasını sağlar.}
\label{fig:henkin}
\end{figure}

\begin{teorem}{Gödel'in tamlık teoremi}{godel-tamlik}
Birinci mertebe mantıkta, her $\Gamma$ cümle kümesi ve her $\varphi$
cümlesi için $\Gamma \models \varphi$ ise $\Gamma \vdash \varphi$.
\end{teorem}

\noindent Bu teoremin şaşırtıcı sonuçları vardır. Tıkızlık birinci mertebe
mantık için de geçerlidir. Sayılabilir bir dildeki her tutarlı kuramın
sayılabilir bir modeli vardır (aşağı Löwenheim--Skolem teoremi): Terim
modeli sayılabilirdir. Gerçel sayılar kuramının bile sayılabilir bir
modeli vardır! Bu ``paradoks'' \ciltiki{5}'te küme kuramıyla birlikte ele
alınacak.

\begin{dikkat}[Dikkat: Tamlık ve eksiklik]
Gödel'in \emph{tamlık} teoremi (1930) ile \emph{eksiklik} teoremleri
(1931) çelişmez, çünkü farklı şeylerden söz ederler. Tamlık, \emph{ispat
sistemi} hakkındadır: Bir kuramın bütün modellerinde doğru olan her cümle
o kuramdan türetilebilir. Eksiklik ise \emph{kuramlar} hakkındadır:
Aritmetiğin, örneğin Peano aritmetiğinin, hem $\varphi$'yi hem
${\sim}\varphi$'yi doğrulayan farklı modelleri olduğu için ikisi de
türetilemez. ``Standart modelde doğru'' ile ``her modelde doğru''
arasındaki fark, \ciltiki{4}'ün konusudur.
\end{dikkat}

% ---------------------------------------------------------------------
\section{Formalizasyondaki tasarım kararları}\label{sec:tamlik-tasarim}

Bu bölümün \Lean{} ispatı birkaç tasarım kararının ürünüdür. Bunlar,
kendi metamantık formalizasyonunuzu yaparken karşılaşacağınız
seçimlerdir.

\begin{table}[htbp]
\centering
\small
\renewcommand{\arraystretch}{1.45}
\rowcolors{2}{tablosatir}{white}
\begin{tabularx}{\linewidth}{>{\raggedright\arraybackslash}p{3.0cm} >{\raggedright\arraybackslash}X >{\raggedright\arraybackslash}X}
\rowcolor{tabloust}
\tablobaslik{Karar} & \tablobaslik{Bizim seçimimiz} & \tablobaslik{Alternatif ve bedeli} \\
Anlambilim değerleri & \lc{Bool}: \leantac{cases} ile satır satır ilerlemek kolay & \lc{Prop}: daha genel, ama RAA ve Kalmár için klasik mantık gerekir \\
Bağlam & \lc{List Formul}: sonlu, \leantac{decide} ile denetlenebilir & \lc{Set Formul}: sonsuz $\Gamma$ mümkün, ama tıkızlık ayrıca ispatlanmalı \\
Tamlık yöntemi & Kalmár: yapıcı, sonlu, ek kütüphane gerektirmez & Lindenbaum: birinci mertebeye genellenir, ama sayma ve seçim gerektirir \\
Atomlar & \lc{Nat} indisli: \lc{enBuyukAtom} ile sınırlanır & Herhangi bir tip: atomların sonlu kümesini ayrıca hesaplamak gerekir \\
Önsav biçimi & \lc{Kalmar Δ b A}: iki yönlü \lc{∧} & \leankw{if}--\leankw{then}: daha kısa ifade, ama her adımda \leankw{if} açılmalı \\
\end{tabularx}
\caption{Tamlık formalizasyonunda yapılan seçimler.}
\label{tab:tamlik-tasarim}
\end{table}

\noindent Bu kararların hiçbiri ``doğru'' ya da ``yanlış'' değildir. Her
biri, ispatın bir bölümünü kolaylaştırırken başka bir bölümünü
zorlaştırır. Formalizasyonda deneyim, büyük ölçüde bu dengeleri önceden
görebilmektir. Michaelis ve Nipkow, Isabelle'de önermeler mantığının dört
farklı ispat sistemini (doğal çıkarım, ardışık hesap, Hilbert sistemi,
çözünürlük) aynı çatı altında formalize etmiş ve bu tür seçimleri
ayrıntılı biçimde tartışmıştır \citep{michaelis2018}.

% ---------------------------------------------------------------------
\begin{arastirmanotu}
\textbf{Birinci mertebe tamlığın formalizasyonları.} Gödel'in tamlık
teoremi birçok ispat asistanında formalize edilmiştir. Lean~3'teki
\emph{Flypitch} projesi, süreklilik hipotezinin bağımsızlığını
ispatlarken birinci mertebe mantığın tamlığını da formalize etti
\citep{han2020flypitch}. Lean~4'teki Foundation kütüphanesi
\citep{foundation} tamlık, tıkızlık ve Gödel'in eksiklik teoremlerini
birlikte içerir \citep{saitou2026}. Bu projelerde en çok emek isteyen
kısım, çoğu zaman tamlık ispatının kendisi değil, \emph{bağlı
değişkenlerin} ve \emph{yerine koymanın} doğru yönetimidir. Bölüm~14'te
bu sorunu de~Bruijn indisleriyle ele alacağız \citep{debruijn1970}.

\textbf{İspat karmaşıklığı.} Cook--Reckhow programı \citep{cookreckhow1979},
$\mathrm{NP} \neq \mathrm{coNP}$'yi ispat sistemleri için alt sınırlar
yoluyla göstermeyi hedefler. SAT çözücülerinin ürettiği çözünürlük
türetimleri bu alanın pratik yüzüdür. \Lean{}'deki \leantac{bv\_decide}
taktiği, bir SAT çözücünün sonucunu \lc{LRAT} sertifikası olarak
denetler ve böylece tamlık teoreminin ``bir türetim vardır'' vaadini
verimli bir biçimde gerçekleştirir (\ciltiki{6}).
\end{arastirmanotu}

\begin{bolumozeti}
\begin{itemize}
  \item Tamlık teoremi, geçerliliğin tersidir: $\Gamma \models A
        \Rightarrow \Gamma \vdash A$. Geçerlilikle birlikte $\vdash$ ve
        $\models$ aynı ilişkidir (\lc{turetilir\_iff\_sonuc}).
  \item Uzun bir ispat, türetilmiş kurallar, Kalmár önsavı, atom eleme
        ve varsayımları içeri taşıma gibi bağımsız önsavlara bölünerek
        yönetilir.
  \item Kalmár önsavı: Doğruluk tablosunun her satırı, işaretli
        atomlardan $A^v$'nin bir türetimidir. Önsav, hem doğru hem yanlış
        durumu birlikte taşıyacak biçimde güçlendirilir.
  \item Atom eleme: Bir atomda farklılaşan iki satır, $p \oveya {\sim}p$
        üzerinden durum ayrımıyla birleşir. $k$ adımda bağlam boşalır.
  \item Tamlık ispatı yapıcıdır ve gerçek bir türetim üretir; ancak
        türetimin boyutu atom sayısında üsteldir.
  \item Sonsuz $\Gamma$ ve birinci mertebe mantık için Lindenbaum önsavı,
        maksimal tutarlı kümeler ve Henkin'in terim modeli kullanılır.
        Tıkızlık ve Löwenheim--Skolem teoremleri bu yöntemin
        sonuçlarıdır.
\end{itemize}
\end{bolumozeti}

\begin{alistirmalar}
\alistirma{1} $\vdash p \oise p \oveya q$. Türetim yazmayın: Doğruluk
tablosunu kontrol edip \lc{tamlik\_bos} kullanın. (\lc{a12\_1})

\alistirma{1} Çift değilleme giriş: $\Delta \vdash A$ ise $\Delta \vdash
{\sim}{\sim}A$. Bu kural RAA gerektiriyor mu? (\lc{a12\_2})

\alistirma{2} $p \ove q \vdash q \ove p$. Bölüm~11'deki \lc{a11\_3}'ü
bu kez \lc{tamlik} ile ispatlayın ve iki ispatı karşılaştırın.
(\lc{a12\_3})

\alistirma{2} Kontrapozisyon, türetilmiş kural olarak: $\Delta \vdash A
\oise B$ ise $\Delta \vdash {\sim}B \oise {\sim}A$. (\lc{a12\_4})

\alistirma{2} $p \oise q \nvdash q \oise p$. (\lc{a12\_5})

\alistirma{1} Kesme kuralı: $\Gamma \vdash A$ ve $A, \Gamma \vdash B$ ise
$\Gamma \vdash B$. Bu kuralın sistemde neden \emph{temel} kural olarak
bulunmadığını tartışın. (\lc{a12\_6})

\alistirma{2} Denklik bağlacı $A \leftrightarrow B$ (\lc{esd A B}) için
Kalmár adımını ispatlayın. \lc{esd}'nin tanımı sayesinde tek satır
yeterli olmalı. (\lc{a12\_7})

\alistirma{3} Tutarlılık ile sağlanabilirlik: $[A] \nvdash \bot \iff A$
sağlanabilirdir. (\lc{a12\_8})

\alistirma{2} (Kâğıt-kalem) $(p \ove q) \oise p$ için
\cref{fig:kalmar-tablo} benzeri dört satırlık tabloyu yazın. Her satırda
hangi türetilmiş kuralların kullanıldığını belirtin.

\alistirma{3} (Proje) Model varlığı teoremini \lc{Classical} kullanmadan
ispatlayın. İpucu: \lc{tumDegerlemeler}'in her değerlemeyi
\lc{enBuyukAtom}'a kadar ``kapsadığını'' gösterin, sonra listeyi
\lc{List.find?} ile tarayın.

\alistirma{3} (Proje) \lc{Turetilir}'in RAA'sız yapıcı versiyonu
\lc{TuretilirY} için Kalmár önsavının hangi adımı başarısız olur? İpucu:
\lc{ucuncu\_hal}.

\alistirma{3} (Araştırma) \lc{tamlik\_bos}'un ürettiği türetimin
boyutunu (kural sayısını) hesaplayan bir fonksiyon yazmak için
\lc{Turetilir}'i \lc{Prop} yerine \lc{Type} değerli yapmak gerekir.
Neden? Kanıt-ilgisizliği (Bölüm~9) burada nasıl bir engel oluşturuyor?
\end{alistirmalar}
